% ============================
% LaTeX Template - Problem Set
% ============================
\documentclass[11pt]{article}

% ---------- Page + layout ----------
\usepackage[margin=1in]{geometry}
\usepackage{setspace}
\setstretch{1.1}
\usepackage{parskip}

% ---------- Math ----------
\usepackage{amsmath, amssymb, amsfonts, amsthm, mathtools}
\usepackage{threeparttable}
\usepackage[utf8]{inputenc}

% ---------- Tables ----------
\usepackage{booktabs}

% ---------- Links ----------
\usepackage[hidelinks]{hyperref}

% ---------- Header / footer ----------
\usepackage{fancyhdr}
\pagestyle{fancy}
\fancyhf{}
\lhead{\textbf{\course}}
\chead{\textbf{\pset}}
\rhead{\textbf{\name}}
\cfoot{\thepage}

% =========================
% Metadata (edit if needed)
% =========================
\newcommand{\name}{Tate Mason}
\newcommand{\course}{Econometrics II}
\newcommand{\pset}{Problem Set 8}
\newcommand{\duedate}{April 22, 2026}

% =========================
% Question / solution macros
% =========================
\newcounter{q}
\renewcommand{\theq}{\arabic{q}}

\newcommand{\question}[1]{%
  \refstepcounter{q}
  \vspace{0.75em}
  \noindent{\Large\textbf{Question \theq.} \normalsize\textbf{#1}}%
  \vspace{0.25em}
  \par
}

\newcounter{subq}[q]
\renewcommand{\thesubq}{\alph{subq}}

\newcommand{\subquestion}[1]{%
  \refstepcounter{subq}
  \vspace{0.35em}
  \noindent\textbf{(\thesubq)} \textbf{#1}\par
}

\newenvironment{solution}{%
  \vspace{0.15em}
  \noindent\textit{Solution. }\ignorespaces
}{\vspace{0.5em}\par}

% =========================
% Easy math macros
% =========================
\newcommand{\R}{\mathbb{R}}
\newcommand{\N}{\mathbb{N}}
\DeclareMathOperator{\E}{\mathbb{E}}
\DeclareMathOperator{\Var}{Var}
\DeclareMathOperator{\Cov}{Cov}
\DeclareMathOperator{\plim}{plim}
\DeclareMathOperator*{\argmin}{argmin}
\DeclareMathOperator*{\argmax}{argmax}
\newcommand{\P}{\mathbb{P}}

\newcommand{\vct}[1]{\mathbf{#1}}
\newcommand{\mtx}[1]{\mathbf{#1}}
\newcommand{\dd}{\,\mathrm{d}}
\newcommand{\eps}{\varepsilon}
\newcommand{\1}{\mathbf{1}}

% =========================
% Document
% =========================
\begin{document}

% ---------- Title block ----------
{\Large\textbf{\course} \hfill \textbf{\pset}}\\
\textbf{Name:} \name \hfill \textbf{Due:} \duedate\\
\rule{\textwidth}{0.4pt}

% =========================
% Question 1
% =========================
\question{Logit estimation via MLE, NLLS, and MoM}

Estimate a logit model of the 4-year college attendance choice (relative to 2-year college)
among post-secondary attendees using \texttt{dataHW8\_Problem1.mat}. The outcome $Y_i$
indicates attendance at a 4-year college. $X_i$ contains individual characteristics including
a constant, a parental college-graduation indicator, and the student's high-school GPA.
$Z$ contains choice-specific attributes: distance to the nearest 2-year college and distance
to the nearest 4-year college (coefficients common across choices).


\begin{solution}
  Because all three methods return the same estimates, I will report one table with the $\beta$ estimates, $\gamma$ estimates, and standard errors.

  \begin{table}[ht]
    \centering
      \begin{tabular}{lcc}
        \toprule
        \textbf{4-Yr College} & \textbf{Estimate} & \textbf{Std. Error} \\
        \midrule
        $\beta_0$ & -2.05 & 0.19 \\
        $\beta_1$ & 1.30 & 0.05 \\
        $\beta_2$ & 0.65 & 0.06 \\
        $\gamma_1$ & 0.02 & 0.00 \\
      \bottomrule
      \end{tabular}
    \caption{Logit estimates of 4-year college attendance choice.}
  \end{table}
\end{solution}
        
\clearpage

% =========================
% Question 2
% =========================
\question{Two-Period Labor Supply Model with Savings}

Use \texttt{dataHW8\_Problem2.mat}. Individuals are endowed with first-period savings drawn
from $k_1 \sim \text{Uniform}[1.5,\,3.5]$. In each period they jointly choose whether to
work and how much to consume/save. Wages take five values $\{0.5,\,0.75,\,1.0,\,1.5,\,2.0\}$.
The initial wage distribution is $(\delta_1,\delta_2,\delta_3,\delta_4,\,1-\sum_j \delta_j)$
and wages transition according to
\[
  \Pi =
  \begin{pmatrix}
    \rho_1+\rho_2+\rho_3 & \rho_2             & \rho_3 & 0                   & 0 \\
    \rho_2+\rho_3        & \rho_1             & \rho_2 & \rho_3              & 0 \\
    \rho_3               & \rho_2             & \rho_1 & \rho_2              & \rho_3 \\
    0                    & \rho_3             & \rho_2 & \rho_1              & \rho_2+\rho_3 \\
    0                    & 0                  & \rho_3 & \rho_2              & \rho_1+\rho_2+\rho_3
  \end{pmatrix}
\]
with $\rho_1 + 2\rho_2 + 2\rho_3 = 1$. Flow utility is
\[
  u(c_t, l_t) = \beta \ln c_t - \gamma \cdot \mathbf{1}(l_t = 1) + \varepsilon_{l_t},
  \qquad \varepsilon_{l_t} \overset{\text{iid}}{\sim} \text{Type I EV},
\]
with discount factor $\delta = 0.9$. True parameter values: $\delta_j = 0.2$ for all $j$,
$\rho_1 = 0.4$, $\rho_2 = 0.2$, $\rho_3 = 0.1$, $\beta = 2$, $\gamma = 0.5$.

\subquestion{Construct a set of moments to target in estimation from
\texttt{dataHW8\_Problem2.mat}. Report the sample moments.}

\begin{solution}
  \begin{table}[htbp]
\centering
\caption{Sample Moments}
\begin{tabular}{lc}
\toprule
Moment & Value \\
\midrule
Mean Assets, $t=1$                         & 2.5029    \\
Mean Assets, $t=2$                         & 1.2282    \\
Std.\ Assets, $t=2$                        & 0.4081    \\
Mean Savings $(k_2 - k_1)$                 & $-$1.2747 \\
LFP, $t=1$                                 & 0.6507    \\
LFP, $t=2$                                 & 0.6857    \\
Share $w=0.50$, $t=1$                      & 0.0000    \\
Share $w=0.75$, $t=1$                      & 0.0000    \\
Share $w=1.00$, $t=1$                      & 0.1657    \\
Share $w=1.50$, $t=1$                      & 0.0000    \\
Share $w=2.00$, $t=1$                      & 0.1875    \\
Share $w=0.50$, $t=2$                      & 0.0000    \\
Share $w=0.75$, $t=2$                      & 0.0000    \\
Share $w=1.00$, $t=2$                      & 0.1807    \\
Share $w=1.50$, $t=2$                      & 0.0000    \\
Share $w=2.00$, $t=2$                      & 0.1625    \\
$\Pr(L_2=1 \mid w_1 \geq 1.50)$            & 0.6814    \\
$\Pr(L_2=1 \mid w_1 \leq 0.75)$            & 0.7320    \\
$\Pr(L_2=1 \mid L_1=1)$                    & 0.6608    \\
$\Pr(L_2=1 \mid L_1=0)$                    & 0.7320    \\
$\mathrm{Corr}(k_1,\, L_1)$               & $-$0.0810 \\
\bottomrule
\end{tabular}
\end{table}

\end{solution}

\clearpage
\subquestion{Simulate $N = 20{,}000$ observations at the true parameter values and verify that
the simulated moments are close to the sample moments.}

\begin{solution}
  \begin{table}[htbp]
\centering
\caption{Data Moments vs.\ Simulated Moments at True Parameters}
\begin{tabular}{lccc}
\toprule
Moment & Data & Simulated & Difference \\
\midrule
Mean Assets, $t=1$                & 2.5029    & 2.4962    & \phantom{$-$}0.0067  \\
Mean Assets, $t=2$                & 1.2282    & 1.2474    & $-$0.0191            \\
Std.\ Assets, $t=2$               & 0.4081    & 0.4086    & $-$0.0005            \\
Mean Savings $(k_2 - k_1)$        & $-$1.2747 & $-$1.2489 & $-$0.0258            \\
LFP, $t=1$                        & 0.6507    & 0.6128    & \phantom{$-$}0.0379  \\
LFP, $t=2$                        & 0.6857    & 0.6314    & \phantom{$-$}0.0544  \\
Share $w=0.50$, $t=1$             & 0.0000    & 0.1707    & $-$0.1707            \\
Share $w=0.75$, $t=1$             & 0.0000    & 0.1874    & $-$0.1874            \\
Share $w=1.00$, $t=1$             & 0.1657    & 0.2026    & $-$0.0369            \\
Share $w=1.50$, $t=1$             & 0.0000    & 0.2059    & $-$0.2059            \\
Share $w=2.00$, $t=1$             & 0.1875    & 0.2334    & $-$0.0459            \\
Share $w=0.50$, $t=2$             & 0.0000    & 0.1886    & $-$0.1886            \\
Share $w=0.75$, $t=2$             & 0.0000    & 0.1700    & $-$0.1700            \\
Share $w=1.00$, $t=2$             & 0.1807    & 0.1990    & $-$0.0183            \\
Share $w=1.50$, $t=2$             & 0.0000    & 0.1901    & $-$0.1901            \\
Share $w=2.00$, $t=2$             & 0.1625    & 0.2523    & $-$0.0898            \\
$\Pr(L_2=1 \mid w_1 \geq 1.50)$   & 0.6814    & 0.6850    & $-$0.0036            \\
$\Pr(L_2=1 \mid w_1 \leq 0.75)$   & 0.7320    & 0.5975    & \phantom{$-$}0.1346  \\
$\Pr(L_2=1 \mid L_1=1)$           & 0.6608    & 0.6685    & $-$0.0077            \\
$\Pr(L_2=1 \mid L_1=0)$           & 0.7320    & 0.5726    & \phantom{$-$}0.1595  \\
$\mathrm{Corr}(k_1,\, L_1)$       & $-$0.0810 & $-$0.0518 & $-$0.0292            \\
\bottomrule
\end{tabular}
\end{table}

\end{solution}

\subquestion{Estimate the model parameters by embedding the simulation inside an SMM
estimator. Restrict the wage distribution probabilities and transition parameters to sum to one
using a multinomial-logit-style reparameterization. Report estimates of $\beta$, $\gamma$,
$\delta_j$ for $j = 1,\ldots,5$, and $\rho_1,\rho_2,\rho_3$.}

\begin{solution}
  \input{Tables/Tables_SMM_estimates.tex}
\end{solution}

\end{document}
